\section{Startup Minimum Viable Safety}

大型平台可以建设专职 trust-and-safety 团队，初创公司却常在用户生成内容刚出现时就面对风险，却缺乏经验、预算和优先级。等待“规模更大再治理”会让数据模型、日志、举报流程和文化全部形成技术债。本节给出 minimum viable safety：不是复制大公司的组织，而是建立最小可升级闭环。

\subsection{Risk inventory、known-content control 与 incident readiness}

创业团队应先列出产品允许谁与谁交互、哪些内容可上传、未成年人是否可见、私信与群组如何传播、哪些第三方模型或存储参与处理。随后选择已知内容匹配、用户举报、moderation queue、account action、audit log 和 emergency escalation 的最低组合，并明确 build、buy 或 specialist partnership。

\lecturefigure{15-startup-maturity.png}{初创公司需要从 risk inventory 走向可升级的 minimum viable safety path。}{本地字幕 35:20--37:20；基于课堂 startup barrier 讨论的概念重绘。}

\begin{knowledgebox}{读图：先让责任和事件有去处}
Risk inventory 回答产品暴露面；minimum controls 提供 known-hash、report 和 basic review；managed program 增加 predictive triage、policy、metrics 与审计；proactive prevention 再加入 threat research、red teaming 与用户保护。第一阶段最重要的不是买齐工具，而是指定 owner、保存必要事件、建立外部报告路径，并能在真实 incident 中找到值班人与决策记录。
\end{knowledgebox}

\begin{importantbox}{Minimum viable safety checklist}
\begin{enumerate}
  \item 一个拥有决策权的 safety owner 与 on-call escalation；
  \item 产品 risk inventory、年龄与交互模型、可接受使用政策；
  \item 用户可见且易用的 report/block/help path；
  \item known-content control、basic queue、case log 与重复事件去重；
  \item 数据最小化、访问权限、retention 与 incident playbook；
  \item 供应商 SLA、故障降级、法定报告和专业机构联系清单。
\end{enumerate}
\end{importantbox}

\begin{warningbox}{“以后再做”会放大迁移成本}
若早期没有稳定 user/account/content ID，后续无法聚合事件；没有 audit log 就无法解释历史处置；没有 age-aware product model 就只能在 UI 完成后硬塞限制；没有 supplier contract 就难以追踪敏感数据。Safety architecture 越晚加入，越可能要求重写 identity、storage、messaging 和 moderation core。
\end{warningbox}

\teachervoice{讲者把 startup 障碍概括为 awareness、money 和 priority，而不是简单责备创始人。课堂的工程含义是提供共享基础设施、清晰成熟度路径和可负担服务，让团队能从最小闭环开始，而不是在“什么都不做”和“自建完整安全部门”之间二选一。}

\subsection{本章小结}

初创公司的安全起点是清楚产品风险、责任人、举报与事件路径，再逐步加入检测、评估和预防。早期集成通常比事故后清理更便宜，也更容易保留正确的数据和权限边界。

